\section{Conséquences pour Aharonov--Bohm, Bell, EPR}
\label{sec:consequences_bell_ab_epr}

Les trois résultats structurels précédents
--- l'angle d'holonomie comme primitive
(Proposition~\ref{prop:theta_p_unique}), le découplage géométrique
(Théorème~\ref{thm:invariance_geom}), et la forme fermée
arithmétique des corrélations
(Corollaire~\ref{cor:closed_form_fr}) --- modifient la lecture
de trois expériences classiques de la physique fondamentale : la
configuration Aharonov--Bohm, les tests de Bell, et la
configuration EPR originelle. Nous examinons brièvement chacune.

\subsection{Aharonov--Bohm : la phase est $\theta_3$, pas $A_\mu$}
\label{ssec:ab_consequence}

Dans la lecture PT, le déphasage interférométrique
$\Delta\varphi$ observé dans une configuration
Aharonov--Bohm \cite{AharonovBohm1959} n'est pas
$(e/\hbar) \oint A \cdot d\ell$ mais $\theta_3$, l'angle
d'holonomie du canal modulaire $p = 3$ associé au flux entouré
par le trajet de la particule. L'identification entre les deux
formules est tautologique en présence du dictionnaire DIC qui
traduit les unités PT en unités SI ; ce qui distingue la lecture
PT n'est pas la valeur numérique du déphasage mais l'identification
du \emph{support physique} du déphasage.

Cette identification a deux conséquences. D'abord, la phase AB
est un invariant strictement topologique du chemin~$\gamma$ :
elle ne dépend que de la classe d'homotopie de $\gamma$ dans le
complémentaire de la région de flux, et est insensible à toute
déformation continue de $\gamma$ dans cette classe. Cette
topologie n'est pas une propriété de la métrique $G(\mu)$ ; elle
est antérieure à la métrique, dans le sens du
Théorème~\ref{thm:invariance_geom}. Ensuite, l'absence de champ
$F = 0$ dans la région intérieure d'un solénoïde idéal n'est plus
un paradoxe : $F_{\mu\nu}$ n'est pas l'objet qui porte
l'information de la phase. C'est $\theta_3$ qui la porte, et
$\theta_3$ est non-nul même là où $F$ s'annule, parce que
$\theta_3$ est défini par la connexion modulaire de PT, qui
n'est elle-même pas une fonction locale du champ.

\subsection{EPR/Bell : corrélations intra-canal, pas inter-canal}
\label{ssec:bell_consequence}

Les tests de Bell \cite{Bell1964,Aspect1982,Hensen2015}
établissent que les corrélations observées entre détecteurs
spatialement séparés excèdent la borne de la
CHSH \cite{CHSH1969} attendue d'une théorie locale à variables
cachées. Dans la lecture standard, ces corrélations sont
interprétées comme la signature d'une non-localité fondamentale du
champ quantique : l'effondrement d'un état EPR sur une mesure à
gauche détermine instantanément l'état corrélé à droite, à toute
distance, et apparemment plus vite que la lumière (mais sans
permettre de communication superluminale, conformément au
no-signalling).

Dans la lecture PT, cette interprétation se renverse. Les deux
mesures effectuées dans un test de Bell ne portent \emph{pas} sur
deux canaux modulaires distincts (auquel cas, par
Théorème~\ref{thm:invariance_geom}, leur MI serait invariante par
géométrie et de magnitude prédite par la formule fermée
arithmétique~\eqref{cor:closed_form_fr}). Elles portent sur le
\emph{même} canal modulaire --- typiquement $p = 2$ pour les
expériences de polarisation/spin, ou $p = 3$ pour les expériences
de couleur de quarks (analogues). Une corrélation intra-canal
n'est pas l'objet du
Théorème~\ref{thm:invariance_geom} : c'est une corrélation
d'ordre supérieur dans la structure CRT, indépendante de toute
notion de séparation spatiale, parce que le canal modulaire
existe avant la construction de l'espace par projection.

Cette lecture s'accorde avec la signature
empirique des tests les plus précis : la corrélation observée
entre deux détecteurs séparés de plusieurs kilomètres
\cite{Hensen2015} est strictement identique à celle observée à
quelques mètres \cite{Aspect1982}, à la précision des mesures
près. Aucune dépendance significative en la distance n'a jamais
été détectée. Cette absence de dépendance est, dans la lecture
PT, l'attendu structurel : ce qui est mesuré n'est pas dans
l'espace, donc la séparation spatiale ne peut pas modifier le
résultat.

\subsection{Dissolution de la non-localité spatiale}
\label{ssec:dissolution_consequence}

La conséquence ontologique des deux sous-sections précédentes est
la \emph{dissolution} (statut [DISS]) de la non-localité spatiale
au sens fort. Le « paradoxe » d'une information apparemment
transmise plus vite que $c$ disparaît parce que les corrélations
ne se propagent pas dans l'espace : elles existent dans la
structure CRT, et l'espace n'est qu'une projection émergente de
cette structure via la métrique $G(\mu)$.

Le statut [DISS] est, dans la nomenclature PT, plus faible que
[THM] mais plus fort que la simple ré-interprétation. Une
dissolution affirme que le problème, tel qu'il était posé dans le
cadre classique, est mal-posé : ses prémisses
implicites --- l'existence d'un espace classique de fond et la
nécessité que toute corrélation observée s'y propage --- ne sont
pas valides en PT. Le problème ne se résout donc pas dans le sens
où l'on aurait trouvé la « vraie » explication classique des
corrélations EPR ; il se dissout dans le sens où les corrélations
sont vues comme des objets qui n'appartenaient jamais au domaine
spatial.

L'horizon expérimental est alors clair : la PT ne prédit aucun
signal qui puisse séparer deux détecteurs spatialement séparés
au-delà de la corrélation arithmétique CRT, et exclut
catégoriquement tout signal qui dépendrait de la métrique
$G(\mu)$ entre les deux détecteurs. Toute observation positive
d'une telle dépendance falsifie simultanément la PT et la
lecture présentée ici. Réciproquement, la confirmation
expérimentale de l'invariance géométrique des corrélations EPR ---
i.e.\ la prédiction QH1a appliquée au régime relativiste --- serait
une validation directe de la lecture PT du débat fondationnel.
